\section{从语言模型到 Agent System}

课堂在 19 分钟处切换到 Agent。这个转折不是再介绍一个模型架构，而是把前半场的 reasoning procedure 放进持续运行的系统：模型要接收目标、拆解任务、调用工具、观察结果、保存状态并决定是否继续。读后半场时应始终区分 \term{foundation model} 与 \term{agent system}；前者提供语言与推理能力，后者才负责行动闭环。

\subsection{Agent 议程：从“为什么”到“怎样控制”}

开场页列出 agent architecture、actions、computer interaction、long-term memory、personalization、multi-agent communication 和 future directions。它给出的教学顺序很重要：先说明为什么一次模型调用不足，再讨论行动接口，最后才进入并行协作与安全治理。

\lecturefigure{slide-28-agent-introduction.jpg}{Agent 半场的课程议程}{Stanford 课堂视频 00:19:51}

\begin{knowledgebox}{读图：七个主题共同构成系统生命周期}
Architecture 决定组件边界，action interface 决定能做什么，memory 与 personalization 决定跨会话状态，multi-agent 决定怎样扩展吞吐和专长，future directions 则处理 reliability、permissions 与 sandboxing。只做其中一个模块，仍不能得到可长期委托的 agent。
\end{knowledgebox}

\teachervoice{Div 先追问“为什么不只训练一个更大的语言模型”。他的回答不是模型能力不重要，而是单次调用没有 chaining、recursion、persistent state 和 environment feedback；真正任务需要系统，不只是 monolith。}

\subsection{为什么需要 Agent：自然语言只是入口}

Building AI Agents slide 提出一个 2023 年课堂 thesis：人通过自然语言表达意图，AI 操作软件与机器。这个愿景减少用户逐个点击界面的负担，但“100 倍速度”与“五年内发生”属于讲者预测，不是课程已经验证的事实。工程上更稳健的表述是：natural-language interface 可以降低任务表达成本，是否提高效率取决于成功率、确认次数和失败恢复成本。

\lecturefigure{slide-29-building-ai-agents.jpg}{Building AI Agents：自然语言到机器操作的愿景}{Stanford 课堂视频 00:20:30}

\begin{importantbox}{Agent 的闭环定义}
令目标为 $g$、内部状态为 $s_t$、环境观察为 $o_t$，agent 在第 $t$ 步执行
\[
p_t=\operatorname{Plan}(g,s_t,o_t),\qquad
a_t=\operatorname{Act}(p_t),\qquad
s_{t+1}=\operatorname{Update}(s_t,o_{t+1}).
\]
$p_t$ 是计划，$a_t$ 是带权限的动作，$o_{t+1}$ 是环境反馈。若系统只生成 $p_t$ 而不执行、观察和更新，它更接近 assistant 或 planner，而不是完整 agent loop。
\end{importantbox}

\begin{warningbox}{界面自动化不会自动消除成本}
Agent 可以减少人工点击，却引入推理延迟、token 成本、错误动作、网站变化、账号权限和审计需求。评价时应比较 end-to-end task success 与总人工介入，而不是只展示一段顺利录像。
\end{warningbox}

\subsection{Agent Stack：模型只是其中一个组件}

架构页把 model 放在中央，周围连接 memory、planning、tools、reflection 与 environment。前半场的 Chain-of-Thought 和 Tree of Thoughts 在这里成为 planning 方法；calculator、calendar 和 code interpreter 成为 tools；短期与长期状态进入 memory；执行后的检查进入 reflection。

\lecturefigure{slide-30-agent-stack.jpg}{包含 memory、planning、tools 与 reflection 的 Agent Stack}{Stanford 课堂视频 00:22:24}

\begin{knowledgebox}{术语消化：Agent Stack 的六个组件}
\begin{tabularx}{\textwidth}{L{0.18\textwidth}>{\raggedright\arraybackslash}X >{\raggedright\arraybackslash}X}
\toprule
组件 & 解决的问题 & 与本讲的关系 \\
\midrule
Model & 理解目标、生成候选计划 & 是 neural compute unit，不独占系统状态 \\
Planning & 把目标转成有依赖关系的步骤 & 可使用 CoT、ToT、decomposition \\
Memory & 保存当前上下文和跨任务事实 & 区分 context 与 persistent store \\
Tools & 执行精确或外部操作 & calculator、browser、code interpreter、apps \\
Reflection & 检查结果并触发修订 & 需要 verifier 或可观察的环境信号 \\
Environment & 返回真实状态变化 & 决定计划是否与现实保持一致 \\
\bottomrule
\end{tabularx}
\end{knowledgebox}

\teachervoice{课堂把构建 agent 类比为构建 computer：LLM 像 CPU，但还要有 RAM、长期存储、I/O、网络、界面和用户配置。这个类比的教学价值，是阻止我们把模型 benchmark 直接等同于系统能力。}

\subsection{MultiOn：课堂原型而非普遍保证}

MultiOn 页介绍讲者当时构建的 browser agent，并用“能够通过自然语言操作互联网”来动机化后续 demo。该页应作为 2023 年课堂原型的来源元数据来读：它说明 demo 来自什么系统，不代表所有网站、账号和任务都具有相同可靠性。本节先建立证据分层，下一章再逐帧分析原型实际展示了什么。

\lecturefigure{slide-31-multion-api.jpg}{MultiOn browser-agent 原型介绍}{Stanford 课堂视频 00:23:03}

\begin{importantbox}{四层证据边界}
\begin{enumerate}
  \item \textbf{课堂画面}：录像确实显示了某次流程的若干状态；
  \item \textbf{讲者陈述}：例如“零人工干预”“拥有购买权限”，是课堂 capability claim；
  \item \textbf{独立验证}：需要可重复任务集、失败日志、版本与评测协议；
  \item \textbf{部署结论}：还需安全、权限、恢复和合规证据。
\end{enumerate}
讲义只把前两层归因给课堂，不把它们自动升级为第三、四层。
\end{importantbox}

\subsection{本章小结}

Agent 不是“会调用工具的大模型”这一句定义能够概括的对象，而是 model、planning、memory、tools、reflection、environment 与 policy 共同组成的闭环。接下来的 live demos 用具体状态展示行动潜力，也暴露不可逆动作和环境约束为何必须进入设计。

